\clearpage
\fancyhead[L]{\small MISCELLANEOUS PAPERS ON ASTRONOMY}
\SourceLeaf{78}%
\pdfbookmark[1]{Calculation of Easter}{osterfest1800}
\begin{center}
{\large CALCULATION OF EASTER.\par}
\medskip\hrule\medskip
{\small Monatliche Correspondenz zur Beförderung der Erd- und Himmels-Kunde,\par
edited by Baron \textsc{von Zach}. August 1800.\par}
\medskip\hrule
\end{center}
\medskip
This paper leaves aside the usual procedure for determining the date of Easter. That procedure can be found in any guide to mathematical chronology and is easy enough in itself once one knows the meaning and use of the customary technical terms \textit{golden number}, \textit{epact}, \textit{Paschal full-moon date}, \textit{solar cycle} and \textit{dominical letter}, and has the necessary tables at hand. The aim here is to give a purely analytical solution to the problem, independent of those auxiliary concepts and based solely on the simplest operations of arithmetic. Hopefully the solution will be welcome to the amateur who is unfamiliar with that method or who may want to determine the date of Easter on the spot, in circumstances where the necessary aids are unavailable or for a year whose calendar cannot be consulted, and will commend itself to the expert by its simplicity and flexibility. The following rules, which anyone who thinks them worth the trouble will easily be able to learn by heart, apply to the two centuries from 1700 to 1899. They can easily be adapted to any other given century by changing their constants as appropriate and adding a minor exception that follows from the arrangement of our calendar and happens to have no occasion to arise during that period.

\SourceLeaf{79}%
\expandafter\def\csname HistoricalNoteText2\endcsname{The reason is that $a$ can take only 19 different values $(0,\ 1,\ 2\ldots18)$, and therefore $d$ can take only as many, among which the value 29 is absent.}
I. Divide the number of the year for which Easter is to be calculated by 19, by 4 and by 7, and call the remainders of these divisions $a$, $b$ and $c$, respectively. If the division is exact, set the corresponding remainder $=0$. Ignore the quotients altogether. The same applies to the following divisions.

II. Next divide $19a+23$ by 30 and call the remainder $d$.

III. Finally divide $2b+4c+6d+3$ or $2b+4c+6d+4$ by 7, according as the given year lies between 1700 and 1799 or between 1800 and 1899, inclusive, and call the remainder $e$.

Easter then falls on the $22+d+e^{\mathrm{th}}$ of March or, if $d+e$ is greater than 9, on the $d+e-9^{\mathrm{th}}$ of April.

\textit{Examples.} For the year 1744, dividing 1744 by 19 gives the remainder $15=a$. The division by 4 is exact, so $b=0$. Dividing by 7 gives the remainder $1=c$. This gives $19a+23=308$, which on division by 30 leaves the remainder $8=d$. Finally $2b+4c+6d+3=55$, divided by 7, gives the remainder $6=e$. Thus Easter falls on the $22+8+6^{\mathrm{th}}$ of March, or the $14-9$, that is, the 5th of April.

For 1800 the values are $a=14$, $b=0$, $c=1$. Then $19a+23=289$, so $d=19$. Next $2b+4c+6d+4=122$, so $e=3$. Easter therefore falls on the $19+3-9$, that is, the 13th of April.

For 1818, $a=13$, $b=2$, $c=5$. Then $19a+23=270$, so $d=0$. Next $2b+4c+6d+4=28$, so $e=0$. Easter therefore falls on the 22nd of March.

In the last example Easter falls on the earliest possible day, for $d$ and $e$ clearly have their smallest possible values here. On the other hand, Easter can never fall later than the $22+29+6^{\mathrm{th}}$ of March, that is, the 26th of April, since $d$ cannot exceed 29 and $e$ cannot exceed 6. In the eighteenth and nineteenth centuries, however, $d=29$ can never occur\HistoricalNote{2}. The latest possible date for Easter during this period is therefore the 25th of April, which occurs when $d=28$ and $e=6$ hold at the same time. These two conditions come together in the years 1734 and 1886. In other centuries $d=29$ could occur, but %
\SourceLeaf{80}%
this is precisely the case in which the exception mentioned above applies, bringing the value of $d$ back down to 28, so that the 25th of April is the latest possible date of all. A fuller treatment of this point would go too far here.

The analysis by which the formula above is found rests, in fact, on principles of \textit{higher arithmetic} for which I cannot yet refer to any publication. The analysis therefore cannot be presented here in all its simplicity. Even so, what follows will suffice to give an idea of the basis of the rules and to show that they are correct.

I. As is well known, the golden number of a year in our reckoning of time is the remainder obtained by adding 1 to the year number and dividing the sum by 19. For exact division, however, the remainder must be set $=19$. It follows readily that $a+1$ is the golden number of the given year.

II. The Paschal full-moon date is the day of the Paschal full moon. In the 18th and 19th centuries, in a year whose golden number is 1, it falls on the 13th of April. Then throughout the full cycle of 19 years, that is, up to and including the year whose golden number is 19, in each year it falls \textit{either} 11 days earlier \textit{or} 19 days later than in the immediately preceding year, according as it fell in \textit{either} April \textit{or} March in that year, as one can readily see from a table of Paschal full-moon dates. Thus it falls on the 2nd of April in the year whose golden number is 2, on the 22nd of March in the following year, in the year whose golden number is 4, on the 10th of April, and so on. It follows that the Paschal full-moon date never falls \textit{before} the 21st of March or \textit{after} the 19th of April. So if, in the year whose golden number is $a+1$, it is taken to fall on the $21+D^{\mathrm{th}}$ of March (counting the days of April as a continuation of March), $D$ always lies between 0 and 29 inclusive. Thus for $a=0$ the value is $D=23$, for $a=1$ it is $D=23-11$, for $a=2$ it is $D=23-2\times11$, for $a=3$ it is $D=23-2\times11+19$, and so on. In general $D=23-11p+19q$, where $p$ and $q$ are determined by the conditions that $p+q=a$ and that $D$ lie between 0 and 29 inclusive. Thus $D=23+19a-30p$, from which one readily concludes that $D$ is the remainder obtained by dividing $23+19a$ by 30. Consequently $D=d$, or the Paschal full-moon date falls on the $21+d^{\mathrm{th}}$ of March.

\SourceLeaf{81}%
III. Easter itself falls only on the \textit{first} Sunday \textit{after} the Paschal full-moon date, hence at least one and at most seven days later, and therefore certainly no earlier than the $22+d^{\mathrm{th}}$ of March. If Easter is taken to fall on the $22+d+E^{\mathrm{th}}$ of March, $E$ lies between 0 and 6 inclusive and must be determined by the condition that this day be a Sunday. This condition can be stated in purely arithmetical terms. The interval between the $22+d+E^{\mathrm{th}}$ of March in the given year and some fixed Sunday must be a number of days divisible by 7 (a whole number of weeks). A particular Sunday must therefore be fixed, and I choose the 21st of March 1700. Now call the number of the given year $A$, and let $i$ be the number of leap years between 1700 and the year $A$, including that year if it is a leap year. Then $i$ is also the number of leap days that have occurred between the 21st of March 1700 and Easter in the year $A$, and the number of \textit{all} days from the 21st of March 1700 to the $22+d+E^{\mathrm{th}}$ of March in the year $A$
\[
=1+d+E+i+365(A-1700)
\]
It is just as easy to see that between 1700 and 1799
\[
i=\tfrac14(A-b-1700)
\]
whereas between 1800 and 1899
\[
i=\tfrac14(A-b-1700)-1
\]
To determine $E$, then, the condition is that
\[
1+d+E+365(A-1700)+\tfrac14(A-b-1700)
\]
or
\[
d+E+365(A-1700)+\tfrac14(A-b-1700)
\]
must be divisible by 7, according as the year falls between 1700 and 1799 or between 1800 and 1899. Thus the result must also be divisible by 7 if a multiple of 7 is added to that number or subtracted from it, or if that number is subtracted from a multiple of 7. To remove the fraction, I first add $\frac74(A-b-1700)$, which is readily seen to be divisible by 7. This gives me
